\documentclass[11pt,a4paper]{article}

\usepackage[margin=2.4cm]{geometry}
\usepackage{amsmath,amssymb}
\usepackage{booktabs}
\usepackage{array}
\usepackage{siunitx}
\usepackage{xcolor}
\usepackage[colorlinks=true,linkcolor=blue!50!black,urlcolor=blue!55!black]{hyperref}
\usepackage{enumitem}
\usepackage{titlesec}

\titleformat{\section}{\normalfont\Large\bfseries\color{blue!45!black}}{\thesection}{1em}{}

\sisetup{
  table-number-alignment = center,
  detect-weight          = true,
}

\newcommand{\dofs}{\ensuremath{\mathrm{dof}}}

\title{Node-level performance of a CVFEM Navier--Stokes multigrid solver\\
       on an NVIDIA GH200 node}
\author{SFEM / CVFEM spike}
\date{\today}

\begin{document}
\maketitle

\begin{abstract}
We report node-level performance of the semi-structured CVFEM Navier--Stokes
solver in SFEM on a single Alps GH200 node (four sockets, \num{72} Neoverse~V2
cores each, \num{288} cores, \(\approx\SI{870}{\giga\byte}\)). At a fixed core
count the rank\,\(\times\)\,thread decomposition is swept from \(1\times288\) to
\(72\times4\); the optimum is \(4\times72\), one rank per socket, at both problem
sizes tested. The decomposition is shown to be iteration-count invariant, so the
wall-clock comparisons are between runs performing identical work. The dominant
cost at every shape is the coarsest-level solve, a dense LU that each rank stores
in full and assembles through an \(n^2\) allreduce; it is a fixed
\(\approx\SI{5.4}{\second}\) floor that does not amortise with problem size and
that bounds the attainable macro-element resolution. We also quantify the cost of
a defect in the coarse hierarchy's mesh packing, which silently tripled the
iteration count at the shipped default settings.
\end{abstract}

\section{Platform and configuration}

All measurements are taken on one Alps GH200 node: four sockets of \num{72}
Neoverse~V2 cores, \num{288} cores in total, \num{36} NUMA domains, and roughly
\SI{870}{\giga\byte} of memory. The toolchain is the \texttt{prgenv-gnu/24.11:v2}
user environment with cray-mpich under \texttt{srun}; threads are bound with
\texttt{OMP\_PROC\_BIND=true} throughout.

The test problem is the lid-driven cavity discretised with the semi-structured
CVFEM Navier--Stokes operator. Two sizes are used:

\begin{center}
\begin{tabular}{lrrrr}
\toprule
case & macro elements & nodes & \dofs & MG levels \\
\midrule
\(N=16\), refine \num{4} & \num{4096} & \num{274625}  & \num{1098500} & 3 \\
\(N=16\), refine \num{8} & \num{4096} & \num{2146689} & \num{8586756} & 4 \\
\bottomrule
\end{tabular}
\end{center}

The solver is FGMRES preconditioned by geometric multigrid with a Vanka
smoother and Galerkin-assembled coarse operators
(\texttt{SFEM\_GMG\_GALERKIN=2}); the coarsest level is solved by a dense LU. All
runs perform one Newton step with a linear tolerance of \num{1e-4}, which is the
regime the solver is used in; reported iteration counts are therefore the linear
iterations of that step.

\section{Rank\,\(\times\)\,thread decomposition at fixed core count}

Table~\ref{tab:shape-small} and Table~\ref{tab:shape-large} sweep the
decomposition at a constant \num{288} cores. In both, the linear iteration count
is identical at every shape, so the wall-clock figures compare runs performing
the same work.

\begin{table}[htbp]
\centering
\caption{Decomposition sweep at \num{1098500}~\dofs\ (three levels).
Linear iterations are \num{7} at every shape.}
\label{tab:shape-small}
\begin{tabular}{l S[table-format=2.2] S[table-format=2.3] S[table-format=1.4]}
\toprule
shape (ranks \(\times\) threads) & {total (\si{\second})} & {coarse factorisation (\si{\second})} & {solve (\si{\second})} \\
\midrule
\textbf{4\,\(\times\)\,72} & \bfseries 9.31 & 5.351 & 1.2901 \\
8\,\(\times\)\,36          & 11.59          & 7.824 & 1.3259 \\
12\,\(\times\)\,24         & 14.66          & 10.864 & 1.3852 \\
18\,\(\times\)\,16         & 20.24          & 16.214 & 1.5299 \\
36\,\(\times\)\,8          & 28.58          & 24.602 & 1.5454 \\
72\,\(\times\)\,4          & 49.99          & 45.778 & 1.7269 \\
\bottomrule
\end{tabular}
\end{table}

\begin{table}[htbp]
\centering
\caption{Decomposition sweep at \num{8586756}~\dofs\ (four levels),
\(\approx\num{29800}\)~\dofs\ per core. Linear iterations are \num{5} at every
shape. The three fat shapes are from a separate job that repeats \(4\times72\);
the two \(4\times72\) measurements, \SI{11.44}{\second} and \SI{11.07}{\second},
bound the run-to-run variation at about \SI{3}{\percent}.}
\label{tab:shape-large}
\begin{tabular}{l S[table-format=2.2] S[table-format=2.3] S[table-format=1.4]}
\toprule
shape (ranks \(\times\) threads) & {total (\si{\second})} & {coarse factorisation (\si{\second})} & {solve (\si{\second})} \\
\midrule
1\,\(\times\)\,288         & 21.50 & 10.864 & 5.5725 \\
2\,\(\times\)\,144         & 16.49 & 6.117  & 4.2616 \\
\textbf{4\,\(\times\)\,72} & \bfseries 11.07 & 5.322 & 1.6595 \\
\textbf{4\,\(\times\)\,72} & \bfseries 11.44 & 5.414 & 1.6693 \\
8\,\(\times\)\,36          & 12.97 & 7.779  & 1.7866 \\
18\,\(\times\)\,16         & 22.61 & 16.945 & 2.4573 \\
72\,\(\times\)\,4          & 50.74 & 45.385 & 2.4001 \\
\bottomrule
\end{tabular}
\end{table}

The optimum is \(4\times72\) at both sizes: one MPI rank per socket, with each
thread team confined to one socket's memory. The curve is monotone on both sides
of it, and the two directions fail for different reasons.

Towards lean ranks, the coarse factorisation dominates: it grows by \(8.6\times\)
from \(4\times72\) to \(72\times4\) while the solve phase moves only
\SI{34}{\percent}. Towards fat ranks, the failure is NUMA. At \(1\times288\) a
single OpenMP team spans all four sockets: the solve phase degrades by
\(3.4\times\) (\SI{1.66}{\second} to \SI{5.57}{\second}), and the factorisation
itself is \emph{worse} than at four ranks (\SI{10.86}{\second} against
\SI{5.32}{\second}), because one dense \texttt{getrf} spread across four sockets
is slower than four independent factorisations each inside a socket. The
elimination of MPI replication and of the \(n^2\) allreduce at one rank does not
compensate.

At saturation the baseline wins on \emph{both} components: even setting the
coarse factorisation aside, the solve phase alone favours \(4\times72\) by
roughly \(1.5\times\) over the lean shapes.

\section{Where the time is spent}\label{sec:time}

The phase breakdown at \(4\times72\) and \num{8586756}~\dofs\ attributes
\SI{5.414}{\second} of \SI{11.44}{\second} --- \SI{47}{\percent} --- to the
coarsest-level factorisation; at \(72\times4\) the same phase is
\SI{89}{\percent} of the run. Raising the problem size eightfold, from
\num{1098500} to \num{8586756}~\dofs, left that cost essentially unchanged
(\SI{5.351}{\second} against \SI{5.414}{\second} at the baseline shape), because
the coarsest level is fixed by the macro mesh (\(17^3\) nodes, \num{19652}~\dofs)
and is untouched by refinement. It is therefore a fixed floor, not an amortisable
overhead.

Table~\ref{tab:coarse} separates the two components of that phase on a
\num{740772}~\dofs\ case with a \num{13500}-\dofs\ coarsest level, in pure MPI.

\begin{table}[htbp]
\centering
\caption{Coarse-solve components at \num{740772}~\dofs, pure MPI (one thread per
rank). \texttt{getrf} is the per-rank dense factorisation; the remainder of the
phase is dominated by the \(n^2\) allreduce that assembles the coarse matrix on
every rank.}
\label{tab:coarse}
\begin{tabular}{S[table-format=3.0] S[table-format=2.1] S[table-format=2.1] S[table-format=2.2]}
\toprule
{ranks} & {\texttt{getrf} (\si{\second})} & {rest of phase (\si{\second})} & {solve (\si{\second})} \\
\midrule
1   & 45.4 & 0.1  & 12.57 \\
8   & 44.9 & 1.9  & 2.97 \\
32  & 47.4 & 1.3  & 1.64 \\
128 & 47.4 & 5.2  & 1.35 \\
288 & 54.9 & 54.9 & 3.41 \\
\bottomrule
\end{tabular}
\end{table}

The two costs have opposite character. The factorisation is replicated, so its
cost per rank is flat and no parallelism is recovered from it. The assembly
allreduce grows with rank count, from \SI{5.2}{\second} at \num{128} ranks to
\SI{54.9}{\second} at \num{288}. The solve phase, by contrast, scales well:
\SI{12.57}{\second} to \SI{1.35}{\second} from \num{1} to \num{128} ranks, a
factor of \num{9.3}.

\paragraph{Resolution ceiling.} The dense coarse solve costs
\(\mathcal{O}(n^3)\) in time and \(\mathcal{O}(n^2)\) in memory, where \(n\) is
the coarse \dofs\ count and scales as \((N+1)^3\) in the macro resolution \(N\)
--- so roughly \(N^9\). Extrapolating from the measured point at \(N=16\)
(\num{19652}~\dofs, \SI{3.09}{\giga\byte} per rank, \SI{5.4}{\second}), \(N=24\)
would need \(\approx\SI{31}{\giga\byte}\) per rank and minutes of factorisation,
and \(N=32\) about \SI{165}{\giga\byte} per rank. This was observed directly: at
\num{288} ranks and a \num{19652}-\dofs\ coarse level the replicated factor
required \(288\times\SI{3.09}{\giga\byte}\approx\SI{890}{\giga\byte}\) and the run
was terminated by the out-of-memory killer. \emph{The coarse solver, not the
domain decomposition, is what bounds the attainable mesh resolution.}

\section{Iteration-count invariance}

The wall-clock comparisons above are only meaningful if the decompositions
perform identical work. Table~\ref{tab:invariance} confirms this over the full
range of rank counts.

\begin{table}[htbp]
\centering
\caption{Linear iterations against rank count, pure MPI. Block-Jacobi is the
structural control: it inverts one \(4\times4\) block per node, so its owned
output depends only on owned input and it cannot be partition dependent. The
four-level and three-level entries at \num{1098500}~\dofs\ are not a
duplication: \(N=8\) at refine \num{8} and \(N=16\) at refine \num{4} coincide in
degree-of-freedom count while differing in hierarchy depth.}
\label{tab:invariance}
\begin{tabular}{lll}
\toprule
case & smoother & linear iterations \\
\midrule
\num{1098500}~\dofs, 4 levels & Vanka        & \num{6} at 1, 8, 32, 64, 128, 288 ranks \\
\num{740772}~\dofs, 3 levels  & Vanka        & \num{7} at 1, 8, 32, 128, 288 ranks \\
\num{1098500}~\dofs, 3 levels & Vanka        & \num{7} at 1, 8, 32, 128 ranks \\
\num{1098500}~\dofs, 3 levels & block-Jacobi & \num{10} at 1, 8, 32, 128 ranks \\
\num{143748}~\dofs, 3 levels  & block-Jacobi & \num{11} at 1, 8, 32, 64 ranks \\
\bottomrule
\end{tabular}
\end{table}

\section{A packing defect in the coarse hierarchy}

The coarse levels of the hierarchy were packed for SIMD execution, and packing
renumbers a mesh's nodes in place. The hierarchy's transfer operators, however,
do not address the coarse vector through the coarse mesh at all: both
\texttt{sshex8\_\allowbreak hierarchical\_\allowbreak prolongation} and
\texttt{sshex8\_\allowbreak hierarchical\_\allowbreak restriction} take only the \emph{fine} element
table and index coarse degrees of freedom by macro-corner identifier. Renumbering
the coarse mesh therefore desynchronised the coarse operator from the transfers
that feed it. The permutation is the identity when a level fits in a single pack,
which is why the defect was invisible below \num{2048} macro elements --- the
shipped default pack size.

Table~\ref{tab:fix} gives the cost at \num{1098500}~\dofs. Two coarse-operator
modes are shown: Galerkin assembly, the default, and rediscretisation.

\begin{table}[htbp]
\centering
\caption{Effect of the coarse-packing defect and its repair at
\num{1098500}~\dofs, serial, \num{72} threads. Pack size \num{2048} is the
shipped default and yields two packs at this mesh size; \num{4096} yields one
pack and \num{0} disables packing. ``cap'' denotes a run that reached the
\num{1000}-iteration limit without converging, so both its iteration count and
its time are lower bounds.}
\label{tab:fix}
\begin{tabular}{llrrrr}
\toprule
& & \multicolumn{2}{c}{before} & \multicolumn{2}{c}{after} \\
\cmidrule(lr){3-4}\cmidrule(lr){5-6}
coarse operator & pack size & iters & time (\si{\second}) & iters & time (\si{\second}) \\
\midrule
Galerkin          & 0 (off)        & 7    & 5.80   & 7  & 5.79 \\
Galerkin          & \textbf{2048}  & \textbf{22} & \textbf{8.97} & \textbf{7} & \textbf{5.82} \\
Galerkin          & 4096           & 7    & 5.65   & 7  & 5.87 \\
rediscretisation  & 0 (off)        & 17   & 13.03  & 17 & 12.85 \\
rediscretisation  & \textbf{2048}  & \textbf{cap} & \textbf{222.65} & \textbf{17} & \textbf{12.97} \\
rediscretisation  & 4096           & 17   & 30.79  & 17 & 12.98 \\
\bottomrule
\end{tabular}
\end{table}

At the shipped default the Galerkin path required three times the iterations and
\SI{55}{\percent} more wall clock; the rediscretisation path failed to converge
at all. The repair is to decline packing on derefined levels, whose numbering is
load-bearing for the transfers, while leaving the fine level's packing --- the
performance-critical one --- untouched. The unpacked configurations are unchanged
to within run-to-run variation, confirming that the change does not perturb the
path it does not touch.

\section{Limitations}

\begin{itemize}
  \item \textbf{Communication overlap is not exercised.} The library's default
    matrix-free path overlaps the halo exchange with computation over element
    subranges. The CVFEM operator does not implement a scoped apply and aborts on
    that path, so every measurement here uses the blocking
    \texttt{SFEM\_PARALLEL\_MF\_LEGACY=1} path. The cost of the communication
    thread at each decomposition shape therefore remains unmeasured.
  \item \textbf{One Newton step.} Iteration counts and timings are for a single
    Newton step at a linear tolerance of \num{1e-4}. A full nonlinear
    continuation has not been swept across decompositions.
  \item \textbf{One coarse-operator mode for the sweeps.} The decomposition
    sweeps use Galerkin-assembled coarse operators. Substituting an iterative
    (BiCGStab) coarse solve removes the factorisation cost but destroys
    iteration-count invariance: it neither converged to its \num{1e-8} tolerance
    nor exhausted its iteration budget, stalling inconsistently and yielding
    outer counts that varied with decomposition and between repeat runs of the
    same configuration. A preconditioned or genuinely distributed direct coarse
    solve, rather than this substitution, is the route to lifting the ceiling in
    Section~\ref{sec:time}.
  \item \textbf{Single node.} No multi-node results are reported.
\end{itemize}

\end{document}
